% Generated draft: requires mathematical and source review.
\documentclass[11pt]{article}
\usepackage[margin=25mm]{geometry}
\usepackage{fontspec}
\setmainfont{Latin Modern Roman}
\usepackage{amsmath,amssymb,mathtools,unicode-math}
\setmathfont{Latin Modern Math}
\usepackage{enumitem,needspace,xurl,hyperref}
\setlength{\parindent}{0pt}
\setlength{\parskip}{5pt}
\setlength{\emergencystretch}{3em}
\newcommand{\N}{\mathbb{N}}
\newcommand{\Z}{\mathbb{Z}}
\newcommand{\Q}{\mathbb{Q}}
\newcommand{\R}{\mathbb{R}}
\newcommand{\C}{\mathbb{C}}
\newcommand{\F}{\mathbb{F}}
\newcommand{\A}{\mathbb{A}}
\newcommand{\PP}{\mathbb{P}}
\newcommand{\E}{\mathbb{E}}
\newcommand{\eps}{\varepsilon}
\newcommand{\dd}{\,\mathrm{d}}
\begin{document}
\section*{UnsolvedMath ID: 11000258}
\textbf{Source problem code:} catalogue entry
\subsection{Definitions and mathematical statement}
Let \(R=\Q(q_1,q_2)\), the rational function field in the source's parameters, interpreted as independent indeterminates. Let \(B_n\) be the \(n\)-strand braid group and retain its previously defined map \(\rho_n:B_n\to H_n(q_1,q_2)\). The target algebra and ``above map'' are not defined in the supplied excerpt.

\emph{Question 1.} Is that map injective over this generic coefficient field? For its intended values of \(n\), the question is
\[
 \rho_n(\beta)=1\ \Longrightarrow\ \beta=1
 \qquad(\beta\in B_n).
\]
The range of \(n\) is not explicitly stated, and no specialization of \(q_1,q_2\) is imposed.

If the missing setup is the standard two-parameter Hecke-algebra convention,
\[
 H_n(q_1,q_2)=R[B_n]/
 \left\langle(\sigma_i-q_1)(\sigma_i-q_2):1\leq i<n\right\rangle,
\]
then \(\rho_n\) is the canonical group map into the multiplicative units of this algebra. This displayed convention is conditional: different parameter or generator conventions must be checked against the source. The desired group injectivity is not injectivity of the group-algebra quotient \(R[B_n]\to H_n\), and the target is not silently replaced by a Temperley--Lieb or another quotient algebra.
\subsection{Short English statement}
Is the source's braid-group map to the generic two-parameter algebra \(H_n(q_1,q_2)\) faithful over \(\Q(q_1,q_2)\)? The algebra and map are omitted, so the usual Hecke interpretation is stated only conditionally.
\subsection{Sources}
\begin{itemize}
\item Local UnsolvedMath registry, permanent ID 11000258, source code catalogue entry, Question 1, PDF page 312. No individual proposer is supplied.
\item Farb, \emph{Problems on Mapping Class Groups and Related Topics}. Supplied URL: \url{https://www.math.uchicago.edu/~farb/papers/mcgbook.pdf}; recorded access 2026-07-29; extraction pdf-font-structure. The volume contains 329 explicitly typeset headings across chapters, with repeated numbering disambiguated by page.
\item The summary mentions Hecke/Temperley--Lieb-type braid representations and an unexplained ``deformation LL'' connection, but gives no precise citation or verified theorem for this map.
\item Supplied status open and OPEN-TRIAGE. The summary cannot verify the exact injectivity question. Current primary-literature status and redistribution terms for the publicly accessible author/publisher copy remain NEEDS\_REVIEW.
\end{itemize}
\end{document}
